\addcontentsline{toc}{subsubsection}{Solution 29. The Conditioned Cubic Squeeze}

\begin{dlsummary}[title={\bfseries Solution 29. The Conditioned Cubic Squeeze},breakable]

\begin{proof}
\textbf{Step 1. Establish oddness and monotonicity.}

On the plane $a+b+c=0$, both conditions apply, so
\[
f(a^3)+f(b^3)+f(c^3)=3f(abc). \tag{1}
\]
Taking $(a,b,c)=(t,-t,0)$ and using $f(0)=0$ gives
$f(-t^3)=-f(t^3)$. Since every real number is a cube, $f$ is odd.

If $x\geq y$, take
$a=\sqrt[3]{x}$, $b=-\sqrt[3]{y}$, and $c=0$. Their sum is
nonnegative, and the first condition gives
\[
f(x)+f(-y)+f(0)\geq3f(0).
\]
Oddness and $f(0)=0$ imply $f(x)\geq f(y)$. Thus $f$ is
non-decreasing.

\medskip
\noindent
\textbf{Step 2. Derive a boundary identity.}

In (1), take $(a,b,c)=(u,v,-u-v)$. Oddness yields
\[
f((u+v)^3)=f(u^3)+f(v^3)+3f(uv(u+v)). \tag{2}
\]
Expand $f((a+b+c)^3)$ by applying (2) first to $(a+b,c)$ and then
to $(a,b)$. Expanding instead first with $(a+c,b)$ gives, after
cancelling the common terms,
\[
f(ab(a+b))+f((a+b)c(a+b+c))
=f(ac(a+c))+f((a+c)b(a+b+c)). \tag{3}
\]

\medskip
\noindent
\textbf{Step 3. Obtain Jensen's equation and affinity.}

Let $x,y>0$ with $x\ne y$, and set
\[
k=\sqrt{\frac yx}\ne1,
\qquad 2k(k^2-1)t^3=x,
\]
where the nonzero real $t$ exists uniquely. Choose
\[
a=(k+1)t,\qquad b=(k-1)t,\qquad c=k(k-1)t.
\]
Direct calculation gives
\[
ab(a+b)=x,\qquad (a+b)c(a+b+c)=y,
\]
and
\[
ac(a+c)=(a+c)b(a+b+c)=\frac{x+y}{2}.
\]
Substitution in (3) proves
\[
f(x)+f(y)=2f\!\left(\frac{x+y}{2}\right). \tag{4}
\]
This also holds when $x=y$, so (4) holds for all positive $x,y$.

Because $f$ is non-decreasing, the midpoint identity implies that $f$
is affine on $(0,\infty)$. Indeed, repeated midpoint interpolation
first gives linear interpolation at dyadic fractions of any interval;
monotonicity then squeezes the value at every intermediate point between
such dyadic approximations. Thus there are constants $m,b$ such that
$f(x)=mx+b$ for all $x>0$.

Substitute this form into (2), for positive $u,v$. The left side is
$m(u+v)^3+b$, while the right side is
$m(u^3+v^3+3uv(u+v))+5b=m(u+v)^3+5b$. Hence $b=0$.
Since $f(1)=1$, we get $m=1$, so $f(x)=x$ for $x>0$. Oddness
and $f(0)=0$ extend this to every real $x$.

\medskip
\noindent
\textbf{Verification.}

For the identity function,
\[
a^3+b^3+c^3-3abc
=\frac12(a+b+c)\bigl((a-b)^2+(b-c)^2+(c-a)^2\bigr).
\]
Its sign agrees with the sign of $a+b+c$, so both required
inequalities hold.
\end{proof}

\end{dlsummary}